\documentclass[11pt,a4paper]{article}
\usepackage[utf8]{inputenc}
\usepackage[italian]{babel}
\usepackage{geometry}
\usepackage{hyperref}
\usepackage{listings}
\usepackage{xcolor}
\usepackage{amsmath}
\usepackage{graphicx}
\usepackage{enumitem}
\usepackage{titlesec}
\usepackage{fancyhdr}

\geometry{margin=2.5cm}

% Configurazione colori per codice
\definecolor{codegreen}{rgb}{0,0.6,0}
\definecolor{codegray}{rgb}{0.5,0.5,0.5}
\definecolor{codepurple}{rgb}{0.58,0,0.82}
\definecolor{backcolour}{rgb}{0.95,0.95,0.92}

\lstset{
    backgroundcolor=\color{backcolour},
    commentstyle=\color{codegreen},
    keywordstyle=\color{magenta},
    numberstyle=\tiny\color{codegray},
    stringstyle=\color{codepurple},
    basicstyle=\ttfamily\footnotesize,
    breakatwhitespace=false,
    breaklines=true,
    captionpos=b,
    keepspaces=true,
    numbers=left,
    numbersep=5pt,
    showspaces=false,
    showstringspaces=false,
    showtabs=false,
    tabsize=2,
    language=JavaScript
}

% Header e footer
\pagestyle{fancy}
\fancyhf{}
\fancyhead[L]{\leftmark}
\fancyhead[R]{PROJECTS \& WORKSPACE - Analisi Tecnica}
\fancyfoot[C]{\thepage}

\title{\textbf{Analisi Tecnica: Componenti PROJECTS e WORKSPACE}}
\author{ExtraTracker - Documentazione Interna}
\date{\today}

\begin{document}

\maketitle
\tableofcontents
\newpage

\section{Introduzione}

Questo documento fornisce un'analisi tecnica dettagliata dei componenti \textbf{PROJECTS} e \textbf{WORKSPACE} dell'applicazione ExtraTracker. L'obiettivo è documentare l'architettura attuale, i flussi di dati, le API, e le interazioni tra i vari layer prima di procedere con una riprogettazione completa.

\subsection{Scopo del Documento}

\begin{itemize}
    \item Documentare l'architettura attuale di PROJECTS e WORKSPACE
    \item Descrivere i modelli di dati e le relazioni
    \item Analizzare i flussi di comunicazione frontend-backend
    \item Identificare le dipendenze e le integrazioni
    \item Fornire una base solida per la riprogettazione
\end{itemize}

\subsection{Struttura del Documento}

Il documento è organizzato in sezioni che coprono:
\begin{enumerate}
    \item \textbf{Architettura Generale}: Panoramica del sistema
    \item \textbf{Componente PROJECTS}: Analisi dettagliata
    \item \textbf{Componente WORKSPACE}: Analisi dettagliata
    \item \textbf{Modelli di Dati}: Schema database e relazioni
    \item \textbf{API e Endpoints}: Contratti di comunicazione
    \item \textbf{Flussi di Dati}: Come i dati fluiscono nel sistema
    \item \textbf{Integrazioni}: Dipendenze tra componenti
\end{enumerate}

\newpage

\section{Architettura Generale}

\subsection{Panoramica del Sistema}

Il sistema ExtraTracker utilizza un'architettura \textbf{full-stack} con separazione tra:
\begin{itemize}
    \item \textbf{Frontend}: React + TypeScript (Vite)
    \item \textbf{Backend}: Node.js + Express + MongoDB
    \item \textbf{Pattern}: Multi-tenancy con isolamento per utente
\end{itemize}

\subsection{Unificazione dei Modelli}

\textbf{IMPORTANTE}: Il sistema ha subito una \textbf{unificazione} dei modelli:
\begin{itemize}
    \item \texttt{WorkProject} (legacy) è stato unificato in \texttt{Project}
    \item \texttt{WorkEntry} (legacy) è stato unificato in \texttt{WorkLog}
    \item Entrambi i componenti (PROJECTS e WORKSPACE) ora condividono lo stesso modello \texttt{Project}
\end{itemize}

\subsection{Separazione delle Responsabilità}

\begin{description}
    \item[\textbf{PROJECTS}] Focus su \textbf{monitoraggio salute progetti}:
    \begin{itemize}
        \item Health monitoring (velocity, projected hours, budget usage)
        \item Metriche e statistiche aggregate
        \item Dashboard di controllo progetti
    \end{itemize}
    
    \item[\textbf{WORKSPACE}] Focus su \textbf{work journal e tracking}:
    \begin{itemize}
        \item Timeline di entries (WorkLog)
        \item Gestione TODO per progetto
        \item Vista cronologica del lavoro
        \item Templates per entries
    \end{itemize}
\end{description}

\newpage

\section{Componente PROJECTS}

\subsection{Panoramica}

Il componente PROJECTS è dedicato al \textbf{monitoraggio della salute dei progetti} con focus su:
\begin{itemize}
    \item Calcolo di metriche avanzate (velocity, projected hours)
    \item Health status (healthy/warning/critical/unknown)
    \item Budget tracking e burn rate
    \item Proiezioni basate su progresso dichiarato
\end{itemize}

\subsection{Struttura Frontend}

\subsubsection{Directory Structure}

\begin{lstlisting}[language=bash]
src/features/projects/
├── components/
│   ├── ProjectFormModal.tsx      # Modal creazione/modifica progetto
│   └── ProjectHealthCard.tsx      # Card con metriche salute
├── context/
│   └── ProjectsContext.tsx       # Context React per stato globale
├── pages/
│   └── ProjectsPage.tsx           # Pagina principale
├── ProjectForm.tsx                # Form progetto (legacy)
└── type.ts                        # TypeScript types
\end{lstlisting}

\subsubsection{Context: ProjectsContext}

Il \texttt{ProjectsContext} gestisce lo stato globale dei progetti:

\begin{lstlisting}
interface ProjectsContextType {
  projects: Project[];              // Lista progetti con metriche
  addProject: (params) => void;    // Crea nuovo progetto
  loading: boolean;                 // Stato caricamento
  refreshProjects: () => void;      // Ricarica progetti
}
\end{lstlisting}

\textbf{Caratteristiche}:
\begin{itemize}
    \item Carica progetti con metriche di salute (\texttt{/api/projects-stats})
    \item Memoizzazione con \texttt{useMemo} e \texttt{useCallback} per performance
    \item Gestione errori con toast notifications
    \item Auto-refresh all'autenticazione
\end{itemize}

\subsubsection{Page: ProjectsPage}

La pagina principale mostra:
\begin{enumerate}
    \item \textbf{Header} con statistiche aggregate:
    \begin{itemize}
        \item Ore totali loggate
        \item Valore generato (earnings)
        \item Progetti critici
        \item Progetti senza stima
    \end{itemize}
    
    \item \textbf{Grid progetti} ordinati per priorità:
    \begin{itemize}
        \item Priorità: critical $\rightarrow$ warning $\rightarrow$ healthy $\rightarrow$ unknown
        \item Ogni progetto mostra \texttt{ProjectHealthCard}
    \end{itemize}
    
    \item \textbf{Modal} per creazione nuovo progetto
\end{enumerate}

\subsubsection{Component: ProjectHealthCard}

Card che visualizza le metriche di salute di un progetto:

\textbf{Informazioni mostrate}:
\begin{itemize}
    \item Nome, codice, tipo (CLIENT/PERSONAL)
    \item Health status badge (healthy/warning/critical/unknown)
    \item Barra progresso consumo ore vs progresso dichiarato
    \item Metriche: ore loggate, stima totale, proiezione finale, guadagno
    \item Velocity message (messaggio dinamico basato su calcoli)
    \item Ultimo log
\end{itemize}

\textbf{Calcolo Health Status}:
\begin{lstlisting}
// Per progetti CLIENT con stima:
if (usagePercent < 75) → 'healthy'
if (usagePercent >= 75 && <= 90) → 'warning'
if (usagePercent > 90) → 'critical'

// Se supera budget → sempre 'critical'
// Se manca stima → 'unknown'
\end{lstlisting}

\subsection{Backend: ProjectService}

\subsubsection{Classe BaseService}

\texttt{ProjectService} estende \texttt{BaseService} che fornisce:
\begin{itemize}
    \item CRUD operations base (find, findById, create, update, delete)
    \item Multi-tenancy automatico (filtro per \texttt{user})
    \item Ricerca testuale
    \item Validazioni e error handling
\end{itemize}

\subsubsection{Metodi Custom}

\begin{description}
    \item[\texttt{findActive(tenantScope, options)}] 
    Trova solo progetti con \texttt{status: 'active'}.
    
    \item[\texttt{findWithEntryCount(tenantScope, options)}]
    \textbf{Cruciale per WORKSPACE}: Aggrega progetti con statistiche:
    \begin{itemize}
        \item \texttt{entriesCount}: Numero di WorkLog associati
        \item \texttt{lastEntryDate}: Data ultimo log
        \item \texttt{totalDuration}: Durata totale in minuti
        \item \texttt{streak}: Giorni consecutivi con entries
    \end{itemize}
    Usa \texttt{\$lookup} MongoDB per aggregare da \texttt{worklogs}.
    
    \item[\texttt{findWithHealth(tenantScope)}]
    Ritorna progetti arricchiti con \texttt{ProjectMetrics}:
    \begin{itemize}
        \item Calcola metriche chiamando \texttt{\_mergeProjectsWithStats()}
        \item Aggiunge campo \texttt{metrics} a ogni progetto
    \end{itemize}
    
    \item[\texttt{getProjectStats(tenantScope)}]
    Statistiche aggregate per dashboard/widget.
\end{description}

\subsubsection{Calcolo Metriche di Salute}

Il metodo \texttt{\_buildHealthSnapshot()} calcola:

\textbf{Input}:
\begin{itemize}
    \item \texttt{project}: Progetto con \texttt{estimatedHours}, \texttt{progress}, \texttt{rate}
    \item \texttt{stats}: Statistiche da WorkLog (ore totali, log count, ultimo log)
\end{itemize}

\textbf{Calcoli}:
\begin{enumerate}
    \item \textbf{Velocity}: Se \texttt{progress > 0} e \texttt{totalHours > 0}:
    \begin{lstlisting}
    velocity = totalHours / (progress / 100)
    // Esempio: 10h per 20% → velocity = 50h per 100%
    projectedHours = velocity * 100
    \end{lstlisting}
    
    \item \textbf{Budget Usage}:
    \begin{lstlisting}
    budgetUsagePercent = (totalHours / estimatedHours) * 100
    \end{lstlisting}
    
    \item \textbf{Health Status}:
    \begin{lstlisting}
    if (projectedHours <= estimatedHours) → 'healthy'
    if (projectedHours <= estimatedHours * 1.2) → 'warning'
    if (projectedHours > estimatedHours * 1.2) → 'critical'
    \end{lstlisting}
    
    \item \textbf{Velocity Message}: Messaggio dinamico basato su calcoli
\end{enumerate}

\subsubsection{Validazioni e Sicurezza}

\begin{itemize}
    \item \textbf{Pre-create}: Verifica unicità codice per utente
    \item \textbf{Pre-delete}: Blocca eliminazione se ci sono WorkLog associati
    \item \textbf{Multi-tenancy}: Tutte le query filtrano per \texttt{user}
    \item \textbf{Rate validation}: Required per CLIENT, opzionale per PERSONAL
\end{itemize}

\subsection{API Endpoints}

\subsubsection{GET /api/projects}

Lista tutti i progetti dell'utente (senza metriche).

\textbf{Response}:
\begin{lstlisting}
{
  success: true,
  data: Project[]
}
\end{lstlisting}

\subsubsection{GET /api/projects-stats}

Lista progetti con metriche di salute (usato da PROJECTS page).

\textbf{Response}:
\begin{lstlisting}
{
  success: true,
  data: (Project & { metrics: ProjectMetrics })[]
}
\end{lstlisting}

\subsubsection{POST /api/projects}

Crea nuovo progetto.

\textbf{Body}:
\begin{lstlisting}
{
  name: string;              // Required
  code: string;              // Required, unique per user
  type: 'CLIENT' | 'PERSONAL'; // Default: 'CLIENT'
  rate?: number;             // Required se CLIENT
  budget?: number;           // Opzionale
  description?: string;
  estimatedHours?: number;
  progress?: number;         // 0-100
  color?: string;            // Default: '#6366f1'
  icon?: string;             // Default: '📂'
  status?: 'active' | 'completed' | 'archived';
}
\end{lstlisting}

\subsubsection{PUT /api/projects/:id}

Aggiorna progetto esistente.

\subsubsection{DELETE /api/projects/:id}

Elimina progetto (solo se non ha WorkLog associati).

\newpage

\section{Componente WORKSPACE}

\subsection{Panoramica}

Il componente WORKSPACE è dedicato al \textbf{Work Journal} con focus su:
\begin{itemize}
    \item Timeline cronologica di entries (WorkLog)
    \item Gestione TODO per progetto
    \item Templates per entries (development, documentation, ticket, etc.)
    \item Vista progetti con statistiche (entries count, streak)
\end{itemize}

\subsection{Struttura Frontend}

\subsubsection{Directory Structure}

\begin{lstlisting}[language=bash]
src/features/workspace/
├── components/
│   ├── EntryCard.tsx              # Card singola entry
│   ├── EntryDetailModal.tsx      # Modal dettaglio entry
│   ├── EntryForm.tsx              # Form creazione/modifica entry
│   ├── EntryTimeline.tsx          # Timeline cronologica
│   ├── ProjectCard.tsx            # Card progetto (con stats)
│   ├── ProjectForm.tsx           # Form progetto
│   ├── ProjectList.tsx           # Lista progetti sidebar
│   ├── TodoForm.tsx               # Form TODO
│   ├── TodoItem.tsx              # Item TODO
│   ├── TodoList.tsx              # Lista TODO
│   ├── TodosWidget.tsx           # Widget TODO
│   └── templates/                # Templates per entries
│       ├── DevelopmentTemplate.tsx
│       ├── DocumentationTemplate.tsx
│       ├── FreeformTemplate.tsx
│       ├── TicketTemplate.tsx
│       └── TemplateSelector.tsx
├── context/
│   └── WorkspaceContext.tsx       # Context React globale
├── pages/
│   ├── WorkspacePage.tsx          # Pagina principale
│   └── ProjectDetailPage.tsx      # Dettaglio progetto
├── services/
│   └── workspaceService.ts        # Service layer API
└── types.ts                        # TypeScript types
\end{lstlisting}

\subsubsection{Context: WorkspaceContext}

Il \texttt{WorkspaceContext} gestisce:
\begin{itemize}
    \item \textbf{Progetti}: Lista con statistiche (entriesCount, streak, etc.)
    \item \textbf{Entries}: Lista WorkLog (timeline)
    \item \textbf{Todos}: Gestione TODO per progetto
\end{itemize}

\textbf{Interface}:
\begin{lstlisting}
interface WorkspaceContextType {
  // Projects
  projects: WorkProject[];
  projectsLoading: boolean;
  projectsError: string | null;
  refreshProjects: () => Promise<void>;
  addProject: (data: CreateWorkProjectDTO) => Promise<void>;
  updateProject: (id: string, data: UpdateWorkProjectDTO) => Promise<void>;
  deleteProject: (id: string) => Promise<void>;
  
  // Entries (WorkLog)
  entries: WorkLog[];
  entriesLoading: boolean;
  entriesError: string | null;
  refreshEntries: () => Promise<void>;
  addEntry: (data: CreateWorkLogDTO) => Promise<void>;
  updateEntry: (id: string, data: UpdateWorkLogDTO) => Promise<void>;
  deleteEntry: (id: string) => Promise<void>;
  
  // Todos
  getTodosByProject: (projectId: string) => Promise<WorkTodo[]>;
}
\end{lstlisting}

\textbf{Caratteristiche}:
\begin{itemize}
    \item Carica progetti con \texttt{includeStats=true} (usa \texttt{findWithEntryCount})
    \item Carica entries tramite \texttt{/api/worklogs/feed}
    \item Gestione ottimizzata con \texttt{useCallback} e \texttt{useMemo}
    \item Auto-refresh all'autenticazione
\end{itemize}

\subsubsection{Page: WorkspacePage}

La pagina principale mostra:

\begin{enumerate}
    \item \textbf{Header} con:
    \begin{itemize}
        \item Titolo "Workspace"
        \item Bottoni: Aggiorna, Nuova Entry, Nuovo Progetto
    \end{itemize}
    
    \item \textbf{Filtri}:
    \begin{itemize}
        \item Ricerca full-text (titolo, descrizione, tag)
        \item Filtro per progetto (dropdown)
        \item Sincronizzazione con URL (\texttt{?projectId=xxx})
    \end{itemize}
    
    \item \textbf{Layout a 2 colonne}:
    \begin{itemize}
        \item \textbf{Sinistra}: \texttt{ProjectList} + \texttt{TodosWidget}
        \item \textbf{Destra}: \texttt{EntryTimeline}
    \end{itemize}
    
    \item \textbf{Modals}:
    \begin{itemize}
        \item \texttt{EntryDetailModal}: Dettaglio entry
        \item \texttt{EntryForm}: Creazione/modifica entry
    \end{itemize}
\end{enumerate}

\textbf{Flusso dati}:
\begin{lstlisting}
1. Carica progetti (con stats) → setProjects()
2. Carica entries (WorkLog feed) → setEntries()
3. Filtra entries in base a:
   - selectedProject (da URL o dropdown)
   - searchQuery (ricerca testuale)
4. Ordina per data (più recente prima)
\end{lstlisting}

\subsubsection{Component: EntryTimeline}

Mostra entries raggruppate per data in ordine cronologico inverso.

\textbf{Caratteristiche}:
\begin{itemize}
    \item Raggruppa per \texttt{date} (YYYY-MM-DD)
    \item Mostra totale minuti per giorno
    \item Ogni entry è un \texttt{EntryCard} cliccabile
    \item Click apre \texttt{EntryDetailModal}
\end{itemize}

\subsubsection{Component: ProjectList}

Lista progetti nella sidebar con:
\begin{itemize}
    \item Nome, icona, colore
    \item Statistiche: entriesCount, streak, lastEntryDate
    \item Selezione progetto (highlight)
    \item Form creazione nuovo progetto
\end{itemize}

\subsubsection{Component: TodosWidget}

Widget TODO che mostra:
\begin{itemize}
    \item TODO del progetto selezionato (o tutti se nessuno selezionato)
    \item Filtri per status (pending, in-progress, completed)
    \item Form creazione TODO
    \item Drag \& drop per riordinare (futuro)
\end{itemize}

\subsection{Backend: WorkspaceService}

\subsubsection{Struttura}

\texttt{workspaceService.js} è un \textbf{wrapper} che espone:
\begin{lstlisting}
module.exports = {
    projects: projectService,  // Usa ProjectService
    todos: workTodoService,    // Usa WorkTodoService
};
\end{lstlisting}

\textbf{NOTA}: Non gestisce direttamente entries perché usa \texttt{workLogService} tramite endpoint unificati.

\subsubsection{WorkLogService}

Il service per WorkLog gestisce:
\begin{itemize}
    \item \textbf{Feed}: Lista cronologica di log (\texttt{getWorkspaceFeed})
    \item \textbf{By Month}: Log di un mese specifico
    \item \textbf{By Project}: Log di un progetto
    \item \textbf{Totals}: Statistiche aggregate per periodo
\end{itemize}

\textbf{Caratteristiche}:
\begin{itemize}
    \item Supporta sia log con orari (Timer) che note senza orari (Journal)
    \item Calcolo automatico \texttt{durationMinutes} se \texttt{startTime/endTime} presenti
    \item Validazione sicurezza: verifica che \texttt{projectId} appartenga allo stesso utente
\end{itemize}

\subsection{API Endpoints}

\subsubsection{GET /api/workspace/projects}

Lista progetti con statistiche (per Workspace).

\textbf{Query params}:
\begin{itemize}
    \item \texttt{includeStats=true}: Include entriesCount, streak, etc.
\end{itemize}

\textbf{Response}:
\begin{lstlisting}
{
  success: true,
  data: (Project & {
    entriesCount: number;
    lastEntryDate: string | null;
    totalDuration: number;
    streak: number;
  })[]
}
\end{lstlisting}

\subsubsection{POST /api/workspace/projects}

Crea progetto (con auto-generazione codice se non fornito).

\subsubsection{GET /api/worklogs/feed}

Feed cronologico di WorkLog (usato da Workspace).

\textbf{Query params}:
\begin{itemize}
    \item \texttt{projectId}: Filtra per progetto
    \item \texttt{date}: Filtra per data specifica
    \item \texttt{startDate/endDate}: Range date
    \item \texttt{tags}: Filtra per tag
    \item \texttt{limit}: Limite risultati
\end{itemize}

\subsubsection{POST /api/worklogs}

Crea nuovo WorkLog (entry).

\textbf{Body}:
\begin{lstlisting}
{
  projectId: string;          // Required
  date: string;               // YYYY-MM-DD, required
  title: string;              // Required
  description?: string;
  tags?: string[];
  mood?: 'high' | 'neutral' | 'low';
  isBillable?: boolean;      // Default: true
  startTime?: string;         // HH:mm (opzionale)
  endTime?: string;          // HH:mm (opzionale)
}
\end{lstlisting}

\textbf{NOTA}: Se \texttt{startTime/endTime} mancano, crea una "nota" (Journal) senza durata.

\subsubsection{GET /api/workspace/todos}

Lista TODO con filtri.

\textbf{Query params}:
\begin{itemize}
    \item \texttt{project}: Filtra per progetto
    \item \texttt{status}: Filtra per status
    \item \texttt{priority}: Filtra per priorità
\end{itemize}

\subsubsection{POST /api/workspace/todos}

Crea nuovo TODO.

\textbf{Body}:
\begin{lstlisting}
{
  project: string;            // Project ID, required
  title: string;              // Required
  description?: string;
  priority?: 'low' | 'medium' | 'high' | 'urgent';
  status?: 'pending' | 'in-progress' | 'completed' | 'cancelled';
  dueDate?: string;           // YYYY-MM-DD
  order?: number;            // Per drag & drop
}
\end{lstlisting}

\newpage

\section{Modelli di Dati}

\subsection{Model: Project}

\textbf{Schema MongoDB}:

\begin{lstlisting}
{
  // Identificazione
  name: String,              // Required, max 100 chars
  code: String,              // Required, unique per user, uppercase
  description: String,       // Opzionale, max 500 chars
  
  // Tipo e tariffa
  type: 'CLIENT' | 'PERSONAL', // Default: 'CLIENT'
  rate: Number,               // Required se CLIENT, min 0
  budget: Number,             // Opzionale, min 0
  
  // UI
  icon: String,               // Default: '📂', max 10 chars
  color: String,              // Hex color, default: '#6366f1'
  
  // Tracking
  estimatedHours: Number,     // Opzionale, min 0
  progress: Number,           // 0-100, default 0
  status: 'active' | 'completed' | 'archived', // Default: 'active'
  
  // Multi-tenancy (aggiunto da plugin)
  user: ObjectId,             // Required, immutable
  
  // Timestamps
  createdAt: Date,
  updatedAt: Date
}
\end{lstlisting}

\textbf{Indici}:
\begin{itemize}
    \item \texttt{\{user: 1, code: 1\}} (unique): Unicità codice per utente
    \item \texttt{\{name: 'text', description: 'text'\}}: Ricerca testuale
\end{itemize}

\textbf{Validazioni}:
\begin{itemize}
    \item Pre-save: Normalizza \texttt{code} in uppercase
    \item Pre-save: Valida \texttt{rate} per CLIENT (required)
    \item Pre-create: Verifica unicità \texttt{code} per utente
\end{itemize}

\subsection{Model: WorkLog}

\textbf{Schema MongoDB}:

\begin{lstlisting}
{
  // Relazioni
  projectId: ObjectId,       // Ref: Project, required
  user: ObjectId,             // Multi-tenancy, required
  
  // Dati base
  date: String,               // YYYY-MM-DD, required
  title: String,             // Required, max 200 chars
  description: String,       // Opzionale, max 5000 chars
  tags: [String],            // Array di tag, max 50 chars each
  
  // Tracking
  mood: 'high' | 'neutral' | 'low',
  isBillable: Boolean,       // Default: true
  
  // Orari (opzionali - per Timer mode)
  startTime: String,         // HH:mm, opzionale
  endTime: String,           // HH:mm, opzionale
  
  // Durata
  durationMinutes: Number,    // Calcolata o manuale, default 0
  isManualDuration: Boolean,  // Flag durata manuale, default false
  
  // Timestamps
  createdAt: Date,
  updatedAt: Date
}
\end{lstlisting}

\textbf{Indici}:
\begin{itemize}
    \item \texttt{\{user: 1, date: -1\}}: Query per mese/periodo
    \item \texttt{\{user: 1, projectId: 1, date: -1\}}: Query per progetto
    \item \texttt{\{title: 'text', description: 'text'\}}: Ricerca testuale
    \item \texttt{\{tags: 1\}}: Ricerca per tag
\end{itemize}

\textbf{Logica Durata}:
\begin{lstlisting}
// Pre-save hook:
if (isManualDuration) {
    // Rispetta valore manuale, non ricalcolare
    return;
}

if (startTime && endTime) {
    // Calcola durationMinutes da orari
    durationMinutes = calculateFromTimes(startTime, endTime);
} else {
    // Se orari mancano, durationMinutes = 0 (nota/journal)
    if (!isModified('durationMinutes')) {
        durationMinutes = 0;
    }
}
\end{lstlisting}

\subsection{Model: WorkTodo}

\textbf{Schema MongoDB}:

\begin{lstlisting}
{
  // Relazioni
  project: ObjectId,          // Ref: Project, required
  user: ObjectId,             // Multi-tenancy, required
  
  // Dati
  title: String,              // Required, max 200 chars
  description: String,        // Opzionale, max 1000 chars
  priority: 'low' | 'medium' | 'high' | 'urgent', // Default: 'medium'
  status: 'pending' | 'in-progress' | 'completed' | 'cancelled', // Default: 'pending'
  
  // Date
  dueDate: String,           // YYYY-MM-DD, opzionale
  completedAt: Date,          // Auto-popolato quando status = 'completed'
  
  // UI
  order: Number,              // Per drag & drop, default 0
  
  // Timestamps
  createdAt: Date,
  updatedAt: Date
}
\end{lstlisting}

\textbf{Indici}:
\begin{itemize}
    \item \texttt{\{user: 1, project: 1, status: 1\}}: Query per progetto e status
    \item \texttt{\{user: 1, dueDate: 1\}}: Query TODO in scadenza
    \item \texttt{\{user: 1, priority: 1, createdAt: -1\}}: Ordinamento per priorità
\end{itemize}

\textbf{Logica}:
\begin{itemize}
    \item Pre-save: Auto-popola \texttt{completedAt} quando \texttt{status = 'completed'}
    \item Pre-save: Rimuove \texttt{completedAt} se status cambia da 'completed'
\end{itemize}

\subsection{Relazioni tra Modelli}

\begin{center}
\begin{tabular}{|l|l|l|}
\hline
\textbf{Modello} & \textbf{Relazione} & \textbf{Tipo} \\
\hline
Project & user & Many-to-One (Multi-tenancy) \\
WorkLog & user & Many-to-One (Multi-tenancy) \\
WorkLog & projectId $\rightarrow$ Project & Many-to-One \\
WorkTodo & user & Many-to-One (Multi-tenancy) \\
WorkTodo & project $\rightarrow$ Project & Many-to-One \\
\hline
\end{tabular}
\end{center}

\textbf{Integrità Referenziale}:
\begin{itemize}
    \item \texttt{WorkLog.projectId} deve riferirsi a un \texttt{Project} dello stesso utente
    \item \texttt{WorkTodo.project} deve riferirsi a un \texttt{Project} dello stesso utente
    \item Validazione implementata nel service layer (non a livello DB)
    \item Eliminazione \texttt{Project} bloccata se ci sono \texttt{WorkLog} associati
\end{itemize}

\newpage

\section{Flussi di Dati}

\subsection{Flusso: Caricamento Progetti in PROJECTS}

\begin{enumerate}
    \item \textbf{Frontend}: \texttt{ProjectsContext.refreshProjects()}
    \item \textbf{API Call}: \texttt{GET /api/projects-stats}
    \item \textbf{Controller}: \texttt{api.js} → \texttt{projectService.getProjectStats()}
    \item \textbf{Service}: \texttt{ProjectService.getProjectStats()}
    \begin{itemize}
        \item Chiama \texttt{\_mergeProjectsWithStats()}
        \item Carica progetti: \texttt{projectService.find()}
        \item Aggrega WorkLog: \texttt{\_aggregateWorklogsByProject()}
        \item Calcola metriche: \texttt{\_buildHealthSnapshot()} per ogni progetto
    \end{itemize}
    \item \textbf{Response}: Array di progetti con campo \texttt{metrics}
    \item \textbf{Frontend}: \texttt{setProjects(data)} → Re-render \texttt{ProjectsPage}
\end{enumerate}

\subsection{Flusso: Creazione Progetto}

\begin{enumerate}
    \item \textbf{Frontend}: \texttt{ProjectsContext.addProject(params)}
    \item \textbf{API Call}: \texttt{POST /api/projects}
    \item \textbf{Controller}: Valida body, chiama \texttt{projectService.create()}
    \item \textbf{Service}: \texttt{ProjectService.create()}
    \begin{itemize}
        \item Validazione pre-create: \texttt{beforeCreate()} (verifica codice unico)
        \item Crea documento MongoDB
        \item Registra activity: \texttt{PROJECT\_CREATED}
    \end{itemize}
    \item \textbf{Response}: Progetto creato
    \item \textbf{Frontend}: \texttt{refreshProjects()} → Toast successo
\end{enumerate}

\subsection{Flusso: Caricamento Workspace}

\begin{enumerate}
    \item \textbf{Frontend}: \texttt{WorkspaceContext} mount
    \item \textbf{Parallelo}:
    \begin{itemize}
        \item \textbf{Progetti}: \texttt{GET /api/workspace/projects?includeStats=true}
        \begin{itemize}
            \item Controller: \texttt{workspaceController.getProjects()}
            \item Service: \texttt{projectService.findWithEntryCount()}
            \item Aggregazione MongoDB con \texttt{\$lookup} su \texttt{worklogs}
            \item Calcolo \texttt{streak} per ogni progetto
        \end{itemize}
        \item \textbf{Entries}: \texttt{GET /api/worklogs/feed}
        \begin{itemize}
            \item Controller: \texttt{api.js} → \texttt{workLogService.getWorkspaceFeed()}
            \item Service: Query WorkLog ordinati per data
        \end{itemize}
    \end{itemize}
    \item \textbf{Frontend}: \texttt{setProjects()} + \texttt{setEntries()}
    \item \textbf{Page}: Filtra e mostra entries in timeline
\end{enumerate}

\subsection{Flusso: Creazione Entry (WorkLog)}

\begin{enumerate}
    \item \textbf{Frontend}: \texttt{WorkspaceContext.addEntry(data)}
    \item \textbf{API Call}: \texttt{POST /api/worklogs}
    \item \textbf{Controller}: Valida body, chiama \texttt{workLogService.create()}
    \item \textbf{Service}: \texttt{WorkLogService.create()}
    \begin{itemize}
        \item Validazione sicurezza: verifica \texttt{projectId} appartiene allo stesso utente
        \item Pre-save hook: Calcola \texttt{durationMinutes} se orari presenti
        \item Crea documento MongoDB
        \item Emit event: \texttt{WORKLOG\_CREATED} (per SSE/notifiche)
    \end{itemize}
    \item \textbf{Response}: WorkLog creato
    \item \textbf{Frontend}: \texttt{setEntries([...prev, created])} → Toast successo
\end{enumerate}

\subsection{Flusso: Calcolo Metriche Salute}

\begin{enumerate}
    \item \textbf{Trigger}: Chiamata a \texttt{getProjectStats()} o \texttt{findWithHealth()}
    \item \textbf{Aggregazione}: \texttt{\_aggregateWorklogsByProject()}
    \begin{itemize}
        \item Raggruppa WorkLog per \texttt{projectId}
        \item Calcola: \texttt{totalMinutes}, \texttt{logCount}, \texttt{lastLog}
    \end{itemize}
    \item \textbf{Merge}: \texttt{\_mergeProjectsWithStats()}
    \begin{itemize}
        \item Combina progetti con statistiche aggregate
        \item Crea mappa \texttt{statsMap} per lookup veloce
    \end{itemize}
    \item \textbf{Calcolo Metriche}: \texttt{\_buildHealthSnapshot()} per ogni progetto
    \begin{itemize}
        \item Calcola \texttt{velocity} e \texttt{projectedHours}
        \item Determina \texttt{healthStatus}
        \item Genera \texttt{velocityMessage}
    \end{itemize}
    \item \textbf{Return}: Array progetti con campo \texttt{metrics}
\end{enumerate}

\newpage

\section{Integrazioni e Dipendenze}

\subsection{Dipendenze Frontend}

\begin{description}
    \item[\textbf{React Router}] Navigazione tra pagine
    \begin{itemize}
        \item \texttt{/projects} → \texttt{ProjectsPage}
        \item \texttt{/workspace} → \texttt{WorkspacePage}
        \item URL params: \texttt{?projectId=xxx} per filtro
    \end{itemize}
    
    \item[\textbf{Shared Services}]
    \begin{itemize}
        \item \texttt{apiClient}: HTTP client congestione errori
        \item \texttt{toast}: Notifiche utente
        \item \texttt{useFormat}: Formattazione (money, hours, date)
    \end{itemize}
    
    \item[\textbf{Auth Context}] Verifica autenticazione per caricamento dati
\end{description}

\subsection{Dipendenze Backend}

\begin{description}
    \item[\textbf{BaseService}] Classe base per CRUD operations
    \begin{itemize}
        \item Multi-tenancy automatico
        \item Ricerca testuale
        \item Error handling standardizzato
    \end{itemize}
    
    \item[\textbf{Multi-Tenancy Plugin}] Isolamento dati per utente
    \begin{itemize}
        \item Aggiunge campo \texttt{user} a tutti i modelli
        \item Filtra automaticamente query per \texttt{user}
    \end{itemize}
    
    \item[\textbf{Activity Service}] Logging attività utente
    \begin{itemize}
        \item \texttt{PROJECT\_CREATED}
        \item \texttt{PROJECT\_COMPLETED}
    \end{itemize}
    
    \item[\textbf{Event Bus}] Eventi per SSE/notifiche real-time
    \begin{itemize}
        \item \texttt{WORKLOG\_CREATED}
        \item Aggiornamenti progetti
    \end{itemize}
\end{description}

\subsection{Comunicazione Frontend-Backend}

\begin{center}
\begin{tabular}{|l|l|l|}
\hline
\textbf{Feature} & \textbf{Endpoint} & \textbf{Service} \\
\hline
PROJECTS - Lista & \texttt{GET /api/projects-stats} & \texttt{projectService.getProjectStats()} \\
PROJECTS - Crea & \texttt{POST /api/projects} & \texttt{projectService.create()} \\
WORKSPACE - Progetti & \texttt{GET /api/workspace/projects?includeStats=true} & \texttt{projectService.findWithEntryCount()} \\
WORKSPACE - Entries & \texttt{GET /api/worklogs/feed} & \texttt{workLogService.getWorkspaceFeed()} \\
WORKSPACE - Crea Entry & \texttt{POST /api/worklogs} & \texttt{workLogService.create()} \\
WORKSPACE - Todos & \texttt{GET /api/workspace/todos} & \texttt{workTodoService.find()} \\
\hline
\end{tabular}
\end{center}

\subsection{Overlap e Ridondanze}

\textbf{Problemi identificati}:
\begin{enumerate}
    \item \textbf{Doppio endpoint progetti}:
    \begin{itemize}
        \item \texttt{/api/projects} (generico)
        \item \texttt{/api/workspace/projects} (specifico workspace)
    \end{itemize}
    Entrambi usano \texttt{ProjectService}, ma con logiche diverse.
    
    \item \textbf{Doppio context}:
    \begin{itemize}
        \item \texttt{ProjectsContext}: Solo per PROJECTS page
        \item \texttt{WorkspaceContext}: Include anche progetti (per workspace)
    \end{itemize}
    Possibile conflitto se entrambi montati.
    
    \item \textbf{Tipi duplicati}:
    \begin{itemize}
        \item \texttt{Project} (projects/type.ts)
        \item \texttt{WorkProject} (workspace/types.ts) - DEPRECATO ma ancora usato
    \end{itemize}
    
    \item \textbf{Normalizzazione dati}:
    \begin{itemize}
        \item Frontend normalizza \texttt{\_id} → \texttt{id}
        \item Logica duplicata in più service
    \end{itemize}
\end{enumerate}

\newpage

\section{Problemi e Limitazioni Attuali}

\subsection{Architetturali}

\begin{enumerate}
    \item \textbf{Separazione non chiara}:
    \begin{itemize}
        \item PROJECTS e WORKSPACE condividono lo stesso modello \texttt{Project}
        \item Ma hanno logiche di business diverse
        \item Endpoint duplicati creano confusione
    \end{itemize}
    
    \item \textbf{Context duplicati}:
    \begin{itemize}
        \item Due context separati per gestire progetti
        \item Possibile inconsistenza se entrambi montati
        \item Nessuna sincronizzazione tra i due
    \end{itemize}
    
    \item \textbf{Service layer inconsistente}:
    \begin{itemize}
        \item \texttt{workspaceService} è solo un wrapper
        \item Logica reale in \texttt{projectService} e \texttt{workLogService}
        \item Difficile capire dove cercare la logica
    \end{itemize}
\end{enumerate}

\subsection{Performance}

\begin{enumerate}
    \item \textbf{Aggregazioni pesanti}:
    \begin{itemize}
        \item \texttt{findWithEntryCount()} fa \texttt{\$lookup} su tutti i worklogs
        \item Calcolo \texttt{streak} richiede query aggiuntive
        \item Nessuna cache o ottimizzazione
    \end{itemize}
    
    \item \textbf{Query ripetute}:
    \begin{itemize}
        \item Metriche ricalcolate ogni volta
        \item Nessun caching lato server
        \item Frontend ricarica tutto anche per piccoli aggiornamenti
    \end{itemize}
\end{enumerate}

\subsection{UX/UI}

\begin{enumerate}
    \item \textbf{Navigazione confusa}:
    \begin{itemize}
        \item Due pagine separate per progetti
        \item Link da PROJECTS a WORKSPACE ma non viceversa
        \item URL params non sempre sincronizzati
    \end{itemize}
    
    \item \textbf{Form duplicati}:
    \begin{itemize}
        \item \texttt{ProjectForm} in projects/
        \item \texttt{ProjectForm} in workspace/components/
        \item Logica simile ma non condivisa
    \end{itemize}
\end{enumerate}

\subsection{Type Safety}

\begin{enumerate}
    \item \textbf{Tipi legacy}:
    \begin{itemize}
        \item \texttt{WorkProject} marcato DEPRECATO ma ancora usato
        \item \texttt{WorkEntry} marcato DEPRECATO ma ancora referenziato
        \item Migrazione incompleta
    \end{itemize}
    
    \item \textbf{Normalizzazione inconsistente}:
    \begin{itemize}
        \item Alcuni service normalizzano \texttt{\_id} → \texttt{id}
        \item Altri no, causando errori runtime
    \end{itemize}
\end{enumerate}

\newpage

\section{Conclusioni e Raccomandazioni}

\subsection{Stato Attuale}

Il sistema attuale funziona ma presenta:
\begin{itemize}
    \item \textbf{Architettura ibrida}: Unificazione parziale dei modelli
    \item \textbf{Logica duplicata}: Endpoint e context separati
    \item \textbf{Performance non ottimale}: Aggregazioni pesanti senza cache
    \item \textbf{Type safety incompleta}: Tipi legacy ancora presenti
\end{itemize}

\subsection{Raccomandazioni per Riprogettazione}

\begin{enumerate}
    \item \textbf{Unificare completamente}:
    \begin{itemize}
        \item Un solo context per progetti (condiviso)
        \item Un solo set di endpoint (con query params per varianti)
        \item Rimuovere tipi legacy
    \end{itemize}
    
    \item \textbf{Separare responsabilità}:
    \begin{itemize}
        \item PROJECTS: Solo visualizzazione metriche/health
        \item WORKSPACE: Solo journal/timeline/todos
        \item Service layer chiaro e dedicato
    \end{itemize}
    
    \item \textbf{Ottimizzare performance}:
    \begin{itemize}
        \item Cache metriche calcolate
        \item Lazy loading per liste lunghe
        \item Aggregazioni ottimizzate con indici
    \end{itemize}
    
    \item \textbf{Migliorare UX}:
    \begin{itemize}
        \item Navigazione unificata
        \item Form condivisi
        \item Sincronizzazione stato tra pagine
    \end{itemize}
    
    \item \textbf{Type safety completa}:
    \begin{itemize}
        \item Rimuovere tutti i tipi DEPRECATED
        \item Normalizzazione consistente
        \item Validazione runtime con Zod o simile
    \end{itemize}
\end{enumerate}

\subsection{Prossimi Passi}

\begin{enumerate}
    \item \textbf{Analisi requisiti}: Definire esattamente cosa deve fare ogni componente
    \item \textbf{Progettazione architettura}: Definire struttura pulita e scalabile
    \item \textbf{Prototipo}: Implementare versione semplificata
    \item \textbf{Migrazione graduale}: Spostare funzionalità senza breaking changes
    \item \textbf{Testing}: Verificare che tutto funzioni correttamente
\end{enumerate}

\vspace{2cm}

\begin{center}
\textit{Documento creato per supportare la riprogettazione dei componenti PROJECTS e WORKSPACE.}
\end{center}

\end{document}
